\documentclass[12pt,a4paper]{article}
\usepackage[utf8]{inputenc}
\usepackage[spanish]{babel}
\usepackage{geometry}
\usepackage{listings}
\usepackage{xcolor}
\usepackage{parskip}
\usepackage{amsmath}
\usepackage{hyperref}

\geometry{margin=2.5cm}

\lstset{
  basicstyle=\ttfamily\small,
  keywordstyle=\color{blue}\bfseries,
  commentstyle=\color{gray},
  breaklines=true, frame=single, language=Java,
  numbers=left, numberstyle=\tiny\color{gray}
}

\title{\textbf{Problema de Rutas --- BFS y UCS (v1)}\\
  \large Documentación Técnica}
\author{Inteligencia Artificial}
\date{}

\begin{document}
\maketitle
\tableofcontents
\newpage

\section{Descripción del Proyecto}
Visualizador de planificación de rutas sobre un grafo de ciudades mexicanas usando \textbf{BFS} (Búsqueda por Amplitud) y \textbf{UCS} (Búsqueda de Costo Uniforme). El usuario selecciona ciudad origen y destino, y el programa encuentra y muestra el camino óptimo.

\section{Estructura del Grafo}

El grafo se representa como una \textbf{lista de adyacencia}. Cada ciudad tiene una lista de vecinos con sus distancias en kilómetros:

\begin{lstlisting}
Map<String, List<Arista>> grafo = new HashMap<>();

grafo.put("Acapulco", Arrays.asList(
    new Arista("Ciudad de Mexico", 400),
    new Arista("Puebla", 480)
));
// ... más ciudades
\end{lstlisting}

Las ciudades incluyen: Acapulco, Ciudad de México, Puebla, Veracruz, Oaxaca, Poza Rica, Huatulco, Querétaro, Pachuca, entre otras.

\section{BFS --- Búsqueda por Amplitud}

BFS explora todos los vecinos de un nodo antes de avanzar al siguiente nivel. Garantiza el camino con el \textbf{mínimo número de saltos} (ciudades intermedias).

\begin{lstlisting}
Nodo_BFS resolverBFS(String origen, String destino) {
    Queue<Nodo_BFS> cola = new LinkedList<>();
    Set<String> visitados = new HashSet<>();
    cola.add(new Nodo_BFS(origen, 0, null));
    visitados.add(origen);

    while (!cola.isEmpty()) {
        Nodo_BFS actual = cola.poll();
        if (actual.ciudad.equals(destino)) return actual;

        for (Arista a : grafo.get(actual.ciudad)) {
            if (!visitados.contains(a.destino)) {
                visitados.add(a.destino);
                cola.add(new Nodo_BFS(a.destino,
                    actual.costo + a.peso, actual));
            }
        }
    }
    return null;
}
\end{lstlisting}

\section{UCS --- Búsqueda de Costo Uniforme}

UCS usa una \textbf{cola de prioridad} ordenada por costo acumulado. Siempre expande el nodo de menor costo total, garantizando la ruta de \textbf{menor distancia}.

\begin{lstlisting}
Nodo_BFS resolverUCS(String origen, String destino) {
    PriorityQueue<Nodo_BFS> cola =
        new PriorityQueue<>(
            Comparator.comparingInt(n -> n.costo));
    Set<String> visitados = new HashSet<>();
    cola.add(new Nodo_BFS(origen, 0, null));

    while (!cola.isEmpty()) {
        Nodo_BFS actual = cola.poll();
        if (visitados.contains(actual.ciudad)) continue;
        visitados.add(actual.ciudad);

        if (actual.ciudad.equals(destino)) return actual;

        for (Arista a : grafo.get(actual.ciudad)) {
            if (!visitados.contains(a.destino)) {
                cola.add(new Nodo_BFS(a.destino,
                    actual.costo + a.peso, actual));
            }
        }
    }
    return null;
}
\end{lstlisting}

\section{Diferencia entre BFS y UCS}

\begin{center}
\begin{tabular}{|l|c|c|}
\hline
\textbf{Característica} & \textbf{BFS} & \textbf{UCS} \\
\hline
Estructura & Cola FIFO & Cola de prioridad \\
Criterio de expansión & Nivel (saltos) & Costo acumulado \\
Solución óptima & En saltos & En distancia total \\
Considera pesos & No & Sí \\
\hline
\end{tabular}
\end{center}

\section{Reconstrucción del Camino}

Cada nodo guarda referencia al nodo padre. La ruta se reconstruye recursivamente:

\begin{lstlisting}
List<String> reconstruirRuta(Nodo_BFS nodo) {
    List<String> ruta = new ArrayList<>();
    while (nodo != null) {
        ruta.add(0, nodo.ciudad);
        nodo = nodo.padre;
    }
    return ruta;
}
\end{lstlisting}

\section{Interfaz Gráfica}

\begin{itemize}
  \item Imagen del grafo de ciudades como fondo visual.
  \item Combos para seleccionar ciudad origen y destino.
  \item Botones \textbf{BFS} y \textbf{UCS} para ejecutar cada algoritmo.
  \item Resultado mostrado: ruta completa, distancia total y número de ciudades.
\end{itemize}

\end{document}
